\documentclass[12pt]{article}
\usepackage{amsmath, amssymb, booktabs}

\title{CSE400 -- Fundamentals of Probability in Computing}
\author{}
\date{}

\begin{document}

\maketitle

\section*{Lecture 18: Marginal PMF, Correlation, Covariance, Joint Gaussian Random Variables}

\hrule
\vspace{0.5cm}

\section{Marginal PMF}

\subsection{Definition}

For discrete random variables $X$ and $Y$, the joint PMF is:

\[
p_{XY}(x,y) = P(X = x, Y = y)
\]

The marginal PMFs are obtained by summing over the other variable:

\begin{align*}
p_X(x) &= \sum_y p_{XY}(x,y) \\
p_Y(y) &= \sum_x p_{XY}(x,y)
\end{align*}

\subsubsection*{Key Idea}

The marginal PMF of one variable is obtained by summing the joint PMF over all possible values of the other variable.

\hrule
\vspace{0.5cm}

\subsection{Marginal PMF from Joint PMF Table}

\begin{itemize}
\item Row sums give $p_X(x)$
\item Column sums give $p_Y(y)$
\end{itemize}

Mathematically:

\begin{align*}
p_X(x_i) &= \sum_j p_{ij} \\
p_Y(y_j) &= \sum_i p_{ij}
\end{align*}

Normalization:

\begin{align*}
\sum_i p_X(x_i) &= 1 \\
\sum_j p_Y(y_j) &= 1
\end{align*}

\hrule
\vspace{0.5cm}

\subsection{Example}

Given joint PMF:

\[
\begin{array}{c|cc}
 & y=0 & y=1 \\
\hline
x=0 & \frac{1}{8} & \frac{1}{4} \\
x=1 & \frac{3}{8} & \frac{1}{4}
\end{array}
\]

\subsubsection*{Step-by-step Solution with Logic}

\textbf{Step 1: Identify rows and columns}

\begin{itemize}
\item Rows correspond to fixed $x$ values
\item Columns correspond to fixed $y$ values
\end{itemize}

\textbf{Step 2: Compute $p_X(x)$ (row-wise sum)}

\begin{align*}
p_X(0) &= \frac{1}{8} + \frac{1}{4} \\
       &= \frac{1}{8} + \frac{2}{8} \\
       &= \frac{3}{8}
\end{align*}

\begin{align*}
p_X(1) &= \frac{3}{8} + \frac{1}{4} \\
       &= \frac{3}{8} + \frac{2}{8} \\
       &= \frac{5}{8}
\end{align*}

Logic: We sum across columns because we eliminate $Y$ to get marginal of $X$.

\textbf{Step 3: Compute $p_Y(y)$ (column-wise sum)}

\begin{align*}
p_Y(0) &= \frac{1}{8} + \frac{3}{8} = \frac{4}{8} = \frac{1}{2} \\
p_Y(1) &= \frac{1}{4} + \frac{1}{4} = \frac{1}{2}
\end{align*}

Logic: We sum across rows because we eliminate $X$ to get marginal of $Y$.

\textbf{Step 4: Verification}

\begin{align*}
p_X(0) + p_X(1) &= \frac{3}{8} + \frac{5}{8} = 1 \\
p_Y(0) + p_Y(1) &= \frac{1}{2} + \frac{1}{2} = 1
\end{align*}

\hrule
\vspace{0.5cm}

\section{Correlation}

\subsection{Definition}

\[
R_{XY} = E[XY]
\]

\subsection{Interpretation}

\begin{itemize}
\item Measures average value of product $XY$
\item Computed about origin (not centered)
\end{itemize}

\subsubsection*{Important Observations}

$R_{XY}$ depends on:

\begin{itemize}
\item Mean values
\item Spread of variables
\item Dependence between variables
\end{itemize}

If $R_{XY} = 0$ $\Rightarrow$ Variables are orthogonal about origin

Limitation: This is not a clean measure of dependence because it mixes mean, spread, and dependence.

\hrule
\vspace{0.5cm}

\section{Covariance}

\subsection{Definition}

\[
\mathrm{Cov}(X,Y) = E[(X - \mu_X)(Y - \mu_Y)]
\]

Where:

\[
\mu_X = E[X], \quad \mu_Y = E[Y]
\]

\subsection{Interpretation}

Covariance measures how deviations from means vary together.

\begin{itemize}
\item $\mathrm{Cov}(X,Y) > 0$ $\rightarrow$ positive covariance (same direction)
\item $\mathrm{Cov}(X,Y) < 0$ $\rightarrow$ negative covariance (opposite direction)
\item $\mathrm{Cov}(X,Y) = 0$ $\rightarrow$ no linear co-variation
\end{itemize}

Important: Magnitude depends on units $\rightarrow$ cannot measure strength.

\subsection{Algebraic Form (Detailed Derivation)}

Start with definition:

\[
\mathrm{Cov}(X,Y) = E[(X - \mu_X)(Y - \mu_Y)]
\]

Step 1: Expand product

\[
= E[XY - X\mu_Y - Y\mu_X + \mu_X\mu_Y]
\]

Step 2: Apply expectation (linearity)

\[
= E[XY] - \mu_YE[X] - \mu_XE[Y] + \mu_X\mu_Y
\]

Step 3: Substitute means

\[
= E[XY] - \mu_Y\mu_X - \mu_X\mu_Y + \mu_X\mu_Y
\]

Simplify:

\[
\mathrm{Cov}(X,Y) = E[XY] - E[X]E[Y]
\]

\subsection{Properties}

\begin{enumerate}
\item Symmetry \\
$\mathrm{Cov}(X,Y) = \mathrm{Cov}(Y,X)$

\item Variance \\
$\mathrm{Cov}(X,X) = \mathrm{Var}(X)$

\item Linearity \\
$\mathrm{Cov}(aX + b, cY + d) = ac \, \mathrm{Cov}(X,Y)$

\item Independence \\
If $X \perp Y \Rightarrow \mathrm{Cov}(X,Y) = 0$
\end{enumerate}

Note: Zero covariance does not imply independence.

\hrule
\vspace{0.5cm}

\subsection{Example 1}

Given: $Y = -2X + 3$

Find $\mathrm{Cov}(X,Y)$

\subsubsection*{Step-by-step Solution with Logic}

\textbf{Step 1:} Use covariance formula

\[
\mathrm{Cov}(X,Y) = E[XY] - E[X]E[Y]
\]

\textbf{Step 2:} Express $XY$ in terms of $X$

\[
XY = X(-2X + 3) = -2X^2 + 3X
\]

Logic: Convert everything into $X$ so expectations can be computed.

\textbf{Step 3:} Compute $E[XY]$

\[
E[XY] = -2E[X^2] + 3E[X]
\]

\textbf{Step 4:} Compute $E[Y]$

\[
E[Y] = -2E[X] + 3
\]

\textbf{Step 5:} Substitute

\[
\mathrm{Cov}(X,Y) = (-2E[X^2] + 3E[X]) - E[X](-2E[X] + 3)
\]

\textbf{Step 6:} Expand

\[
= -2E[X^2] + 3E[X] + 2(E[X])^2 - 3E[X]
\]

\[
= -2E[X^2] + 2(E[X])^2
\]

\textbf{Step 7:} Use variance identity

\[
\mathrm{Var}(X) = E[X^2] - (E[X])^2
\]

Thus:

\[
\mathrm{Cov}(X,Y) = -2 \, \mathrm{Var}(X)
\]

\hrule
\vspace{0.5cm}

\section{Pearson’s Correlation Coefficient}

\subsection{Definition}

\[
\rho_{XY} = \frac{\mathrm{Cov}(X,Y)}{\sigma_X \sigma_Y}
\]

Where:

\[
\sigma_X = \sqrt{\mathrm{Var}(X)}, \quad \sigma_Y = \sqrt{\mathrm{Var}(Y)}
\]

\subsection{Properties}

\begin{itemize}
\item Measures direction and strength of linear relationship
\item Normalized covariance
\item Unit-free
\end{itemize}

Range:

\[
-1 \leq \rho \leq 1
\]

Interpretation:

\begin{itemize}
\item $\rho > 0$ $\rightarrow$ positive linear trend
\item $\rho < 0$ $\rightarrow$ negative linear trend
\item $\rho = 0$ $\rightarrow$ no linear correlation
\end{itemize}

\hrule
\vspace{0.5cm}

\section{Example 2}

Given:

\[
f_{XY}(x,y) = \frac{xy}{96}, \quad 0 < x < 4, \; 1 < y < 5
\]

\textbf{Step 1: Compute $E[X]$}

\[
E[X] = \frac{1}{8} \int_0^4 x^2 dx = \frac{8}{3}
\]

\textbf{Step 2: Compute $E[Y]$}

\[
E[Y] = \frac{1}{12} \int_1^5 y^2 dy = \frac{31}{9}
\]

\textbf{Step 3: Compute $E[XY]$}

\[
E[XY] = E[X]E[Y] = \frac{248}{27}
\]

\textbf{Step 4: Compute $E[2X + 3Y]$}

\[
E[2X + 3Y] = \frac{47}{3}
\]

\textbf{Step 5: Compute Var(X)}

\[
\mathrm{Var}(X) = \frac{8}{9}
\]

\textbf{Step 6: Compute Var(Y)}

\[
\mathrm{Var}(Y) = \frac{92}{81}
\]

\textbf{Step 7: Compute Cov(X,Y)}

\[
\mathrm{Cov}(X,Y) = 0
\]

\textbf{Step 8: Compute Correlation}

\[
\rho = 0
\]

\hrule
\vspace{0.5cm}

\section{Jointly Gaussian Random Variables}

\subsection{Definition}

A pair $(X,Y)$ is jointly Gaussian if:

\begin{itemize}
\item Joint density has bivariate Gaussian form
\item Marginals are Gaussian
\end{itemize}

\subsection{Joint Gaussian PDF}

\[
f_{XY}(x,y) =
\frac{1}{2\pi\sigma_X\sigma_Y\sqrt{1-\rho^2}}
\exp\left[
-\frac{1}{2(1-\rho^2)}
\left(
\left(\frac{x-\mu_X}{\sigma_X}\right)^2
+
\left(\frac{y-\mu_Y}{\sigma_Y}\right)^2
-2\rho
\left(\frac{x-\mu_X}{\sigma_X}\right)
\left(\frac{y-\mu_Y}{\sigma_Y}\right)
\right)
\right]
\]

\subsection{Interpretation}

\begin{itemize}
\item Used in signal processing
\item Used in communication systems
\item Models noisy measurements
\end{itemize}

Example: CO$_2$ concentration and temperature

If $\rho > 0$ $\rightarrow$ higher CO$_2$ corresponds to higher temperature

\hrule
\vspace{0.5cm}

\begin{center}
\textbf{End of Lecture 18 Scribe}
\end{center}

\end{document}